% Transcription — fonds Alexandre Grothendieck, shelfmark 43, batch 2 (pages 21-40).
% Established from the facsimile archives/batches/43/batch-02.pdf (PDF page k =
% archivists' page 20+k, no cover sheet), cut from a copy of 43.pdf fetched
% through the project's own relay
% (https://grothendieck-relay.onrender.com/source/43.pdf); tiles in
% archives/tiles/43/p21-p40. Per the skill transcribe-grothendieck. The folder
% is 207 pages in eleven batches, transcribed in parallel. Batch 1 (committed)
% and batches 9-11 were read for the folder's conventions; pages 21-22 continue
% the geometric class-field run that ends batch 1 (G-divisors, generalised
% Jacobians J_{O_D/k}, J_{K/k}), and the notation follows batch 1 there.
%
% Pass: Opus 5.5 (claude-opus-5-5), 2026-09-26 — first pass, unchecked against the pages by a human.
%
% Dating: the folder has its own inventory date, carried verbatim in \dating{}
% with no group sentence (the skill adds the group « Jacobiennes » (43-44)
% only for undated folders). No page of this batch carries a date.
%
% Copies. No page of this batch is evidently a photocopy: all are ink (blue,
% blue-black, brown on p. 24, one red marginal on p. 36) with the pressure and
% colour variation of originals; covers 23 and 35 are chamois card.
%
% Pages skipped: 38, a blank cream sheet. Pages 23 and 35 are covers inscribed
% in his hand and take no \page marker; their words become section headings,
% with a \note: p. 23 « Théorème de Rosenlicht », p. 35 « Adèles des
% préschémas ». Page 37 carries only three lines (the end of p. 36) and is
% transcribed.
%
% Pages summarised: none. Page 30 is fast prose with almost no formula and is
% transcribed with a heavy apparatus; the diagrams of pp. 21, 22, 26, 28, 34
% are redrawn with their layout simplified and said in \note{}s. Page 40 is
% written upside down relative to the archivists' number and is read turned.
%
% His pagination: « 1 » at the top left of p. 24, heading the Rosenlicht run
% (pp. 24-34); no further number is visible on pp. 25-34. Numbered runs of
% his: sections (I) « Existence des Picard pour les courbes » (p. 24) and
% (II) (p. 25); hypotheses a)-e) (p. 24); Théorème (Th. de Rosenlicht) and
% Corollaires 1, 2, 3 (pp. 27-29). No internal page reference.
%
% What the batch is about. Pp. 21-22 (end of batch 1's run): exact sequences
% for a commutative algebraic group G on a complete X over k: G(K)/G(C) ->
% Div^0_G -> H^1(X,G)^0, prod J_{O_Z/k} (codim 2) -> prod J^0_{O_D/k} -> J^0_{K/k}
% -> script-J^0_{X/k}, G-divisors as Hom's, (G-Div)/(G-Div)^0 in
% Hom(H_1(J_*), G), Néron-Severi, « A-t-on égalité ?? », « Quel était
% l'accouplement ... de Tate ??? »; p. 22: G-divisors on X x Y as
% correspondences, J*_{X/k} -> Pic*_G(Y), and a diagram from the Albanese
% sequence (0-cycles of degree 0) to 0 -> G(K)/G(Y) -> Div^0(Y,G) ->
% Pic^0(Y,G) -> 0. Pp. 24-34, under the cover « Théorème de Rosenlicht »: (I)
% existence of the relative Picard scheme Pic_{X/Y} of a flat proper family of
% curves (a)-e)), smooth, with augmentation to Z_{Y'/Y} and trace to Z_Y; the
% morphism X - S -> Pic^1_{X/S}; (II) for X proper smooth with curves as
% fibres and S finite flat, the contractions X_n = X amalg_{S_n} Y and the
% pro-group Pic_{S,X/Y} (generalised Jacobian with modulus S), the universal
% map phi_S : X - S -> Pic^1_{S,X/Y}; Théorème (Rosenlicht): Hom of
% Z-augmented group schemes (Pic_S, P) = Hom_Y(X - S, P^1); Cor. 1 (Hom
% into G); Cor. 2 (phi*_S fully faithful from extensions of Pic_S by G to
% G-torsors on X - S), with its proof; Cor. 3 (Ext^1(Pic_S, G) -> H^1(X-S,G)
% injective, « il semble » bijective when S -> Y is surjective); p. 30,
% reduction of surjectivity to formal completions and residue fields;
% pp. 31-32, an exact sequence 0 -> Ext^1(Pic_{X/Y},G) -> H^1(X,G) ->
% Hom(G_m,G), an element zeta = eta o xi in Ext^2(Pic_{X/Y}, G_m) built from
% 0 -> N_S -> Pic_S -> Pic -> 0, and the « plausible » sequences 0 -> Ext^i
% -> H^i -> Ext^{i-1}(G_m,G) -> 0 when zeta = 0; p. 33, the limit
% Pic_infinity, H^i(R_{X/Y},G) = Ext^i(Pic_infinity, G), and the long exact
% sequence from 0 -> Pic^0_S -> Pic_S -> Z -> 0; p. 34 (« Détailler »),
% autoduality of Jacobians, local-to-global spectral sequences, and the
% plausible H^*(X_S/Y,G) = Ext^*_S(Pic_S,G). Pp. 36-37, under the cover
% « Adèles des préschémas »: the ring scheme S_infinity (A[[t]]** with
% convolution), Frobenius F and Verschiebung V (FV = VF = p), and the
% morphism psi : S_infinity -> W_infinity to Witt vectors, with Teichmüller
% representatives as G_m -> W_infinity. Pp. 39-40: Galois descent for the
% idele class group J_C of a curve (Gamma, pi), Tate cohomology of K* and
% K_x* zero, C*(Gamma, J_C) a resolution, a spectral sequence H^p(H^q(pi,
% ...)) => H*(pi, J_C'(k-bar)) with Lang's theorem, and the cochain
% differential (delta phi)(x) = phi(gamma x) - phi(x).
%
% Notation fixed once for the batch: his hooked P is \mathrm{Pic} throughout
% (Pic_{X/Y}, Pic_{S,X/Y}, Pic^1, Pic^0, Pic_infinity); \mathcal{D}iv for his
% curly Div, as batch 1 writes it; J straight and a looped J set \mathcal{J}
% (p. 21-22), kept distinct; \mathfrak{J} for the ideal sheaf (pp. 25-26);
% \mathbb{Z}_{Y'/Y}, \mathbb{G}_m; \mathbb{S}_\infty and \mathbb{W}_\infty for
% his barred / double S and W (p. 36); \underline and double underline on
% H, Ext where he underlines (sheaf versions); \hat H for Tate cohomology.
%
% Readings to check first: the marginal notes of pp. 21, 24, 27, 28, 31;
% « simple » (p. 24, underlined) and « définissant Y » (p. 25, where S is
% expected); the long cancelled complement on p. 28; the whole prose of
% pp. 30, 31, 33; « Détailler » (p. 34); the name of the global theorem on
% p. 39; the circled word over « G-diviseur » on p. 22.

\documentclass[11pt,a4paper]{article}
% Le préambule commun à toutes les transcriptions du fonds Grothendieck.
%
% Il tient en une idée : **l'incertitude se compose, elle ne se résout pas.**
% Une transcription qui lisse un mot douteux a détruit la seule chose pour
% laquelle on la faisait — un lecteur doit pouvoir savoir, à chaque endroit,
% ce qui était sur le feuillet et ce qui a été ajouté. Les six macros qui
% suivent sont donc l'appareil critique, et non de la décoration.
%
% Les mêmes commandes sont reconnues par `scripts/render.mjs`, qui produit la
% vue de lecture du site. Une macro ajoutée ici doit l'être aussi là-bas,
% faute de quoi le rendu échouera bruyamment — ce qui est voulu : un rendu
% silencieusement faux est pire qu'un rendu absent.

\ProvidesPackage{grothendieck}[2026/08/08 Transcriptions du fonds Grothendieck]

\usepackage{amsmath,amssymb,amsthm}
\usepackage{mathrsfs}
\usepackage{stmaryrd}
% Les diagrammes commutatifs sont partout dans ce fonds : c'est la langue
% même de la géométrie algébrique de Grothendieck, et une transcription qui
% les rendrait en prose ne transcrirait rien.
\usepackage{tikz-cd}
% Français par défaut — c'est la langue des feuillets — mais l'anglais est
% chargé aussi : la traduction et le résumé sont des documents anglais, et la
% typographie française (espaces avant les deux-points) y serait une faute.
% Ces documents basculent par \selectlanguage{english} en tête de corps.
\usepackage[english,french]{babel}
\usepackage{fontspec}
\usepackage{geometry}
\geometry{a4paper,margin=25mm,marginparwidth=28mm}
\usepackage{marginnote}
\usepackage{xcolor}
% ulem, not soul. `soul`'s \st and \ul are built on a character-by-character
% scan that XeLaTeX's font machinery breaks: the document fails with an
% undefined control sequence rather than degrading. ulem does the same two jobs
% and is Unicode-engine safe. `normalem` stops it hijacking \emph, which the
% transcriptions use for Grothendieck's own emphasis.
\usepackage[normalem]{ulem}
\usepackage{hyperref}
\hypersetup{colorlinks=true,linkcolor=black,urlcolor=black,pdfborder={0 0 0}}

% --- Métadonnées du lot -----------------------------------------------------
% Reprises telles quelles par le script de rendu, qui les affiche en tête de
% la vue de lecture. Elles fixent ce que le fichier prétend couvrir : sans
% elles, une transcription flotte sans point d'attache dans un fonds de seize
% mille feuillets.

\newcommand{\folder}[1]{\gdef\grothFolder{#1}}
\newcommand{\batch}[1]{\gdef\grothBatch{#1}}
\newcommand{\leaves}[2]{\gdef\grothFirst{#1}\gdef\grothLast{#2}}
\newcommand{\foldertitle}[1]{\gdef\grothFoldertitle{#1}}
\gdef\grothFolder{?}\gdef\grothBatch{?}\gdef\grothFirst{?}\gdef\grothLast{?}\gdef\grothFoldertitle{}

% --- Le résumé --------------------------------------------------------------
%
% Il ouvre la lecture modernisée, sous le titre « Résumé », et il est composé
% en petit. Ce n'est pas un ornement : c'est le seul passage du document qui
% ne soit pas des mathématiques, et un lecteur doit pouvoir le reconnaître
% avant de l'avoir lu — puis le sauter s'il connaît déjà le sujet, ou s'y
% arrêter s'il ne le connaît pas. Le filet qui le suit marque la reprise du
% corps.

\newenvironment{resume}
  {\small\setlength{\parskip}{0.45em}}
  {\par\medskip\hrule\medskip}

% --- Mots-clés ---------------------------------------------------------------
%
% En anglais, en clôture du résumé de la lecture modernisée : le vocabulaire
% moderne sous lequel chercher ce que les feuillets construisent. Le manifeste
% du site (`npm run manifest`) les extrait pour étiqueter la cote entière —
% cette ligne est la source unique des tags, il n'y a pas de fichier à côté.
\newcommand{\keywords}[1]{\par\smallskip\noindent
  {\footnotesize\textsc{Keywords} --- \textit{#1}}\par}

% --- L'appareil critique ----------------------------------------------------

% Le numéro de feuillet, en marge. C'est l'ancre qui permet de régler un
% désaccord sur une formule en regardant *un* feuillet plutôt qu'une chemise
% de six cents. Le site s'en sert aussi pour tourner les pages du fac-similé
% quand on fait défiler la transcription.
\newcommand{\leaf}[1]{%
  \marginnote{\footnotesize\color{gray!70}\textsf{#1}}%
  \label{leaf:#1}\ignorespaces}

% Illisible : on ne devine pas. Un mot inventé qui se lit comme les autres est
% le pire résultat possible.
\newcommand{\ill}{\textcolor{red!60!black}{[\,\ldots\,]}}

% Lecture douteuse : le mot est donné, et signalé comme incertain.
\newcommand{\uncertain}[1]{\uline{#1}}

% Ajout de l'éditeur — une lettre manquante, un mot sous-entendu.
\newcommand{\add}[1]{\textcolor{blue!50!black}{[#1]}}

% Biffé par Grothendieck lui-même. Ce qu'il a rayé fait partie du document :
% une rature dit souvent où la pensée a changé de direction.
\newcommand{\struck}[1]{\sout{#1}}

% Note du transcripteur, sur l'état matériel du feuillet.
\newcommand{\note}[1]{\par\smallskip{\small\color{gray!60!black}\textit{[#1]}}\par\smallskip}

% Note marginale de Grothendieck, distincte du corps du texte.
\newcommand{\marginal}[1]{\par\smallskip\noindent\hspace{1em}%
  {\small\color{gray!60!black}#1}\par\smallskip}

% --- Titre ------------------------------------------------------------------

\AtBeginDocument{%
  \begin{center}
    {\large\bfseries Fonds Alexandre Grothendieck --- cote n\textsuperscript{o}~\grothFolder}\\[2pt]
    {\itshape\grothFoldertitle}\\[4pt]
    {\small feuillets \grothFirst--\grothLast\ (lot \grothBatch)}\\[2pt]
    {\footnotesize\color{gray!60!black}Transcription établie d'après le fac-similé numérisé
     par l'Université de Montpellier.\\
     Les lectures incertaines sont soulignées, les illisibles notées [\,\ldots\,],
     les ajouts entre crochets.}
  \end{center}
  \vspace{1em}}

\endinput


\folder{43}
\batch{2}
\pages{21}{40}
\dating{[à partir de 1961-vers 1968]}
\watermark{Édition de démonstration}
\foldertitle{Dualité des groupes. Jacobiennes généralisées etc : notes et copies de notes manuscrites (s.d.)}

\begin{document}

\page{21}
\note{page de suites exactes et de schémas, dans la continuité des pages 18
et 19 (courbe $X$, groupe $G$, jacobiennes généralisées) ; les formules sont
données dans l'ordre de la page}
\[
0 \to G(K)/G(C) \longrightarrow \mathcal{D}iv^{0}_{G} \ \Big|\ H^1(X, G)^{0} \to 0
\]
\note{le $H$ est écrit sur une lettre antérieure ; un trait vertical coupe
cette ligne et la suivante ; entre les deux lignes, deux traits obliques en
croix, dont l'un, repassé, va de $\prod J^{0}_{\mathcal{O}_D/k}$ vers
$\mathcal{D}iv^{0}_{G}$}
\[
\prod_{Z\ \text{de codim}\ 2} J_{\mathcal{O}_Z/k} \to \underbrace{\prod J^{0}_{\mathcal{O}_D/k}}_{?} \xrightarrow{\ u\ } J^{0}_{K/k} \ \Big|\ \longrightarrow \mathcal{J}^{0}_{X/k} \to 0
\qquad \text{\emph{$X$ complète}}
\]
\note{$\mathcal{J}$ rend un $J$ bouclé, distinct du $J$ droit des autres
termes, et gardé ainsi}
\struck{$\mathcal{D}iv_G \quad G(K)/G(C)$}

\struck{$\mathcal{D}iv \subset \mathcal{D}iv$}
\[
(G\text{-}\mathcal{D}iv) = \mathrm{Hom}\Bigl(\bigl(\prod \mathcal{J}^{0}_{\mathcal{O}_D/k}\bigr) / \mathrm{Im}\bigl(\prod_{Z\ \text{codim}\ 2} J_{\mathcal{O}_Z/k}\bigr)\Bigr)
\]
\[
(G\text{-}\mathcal{D}iv)^{0} \underset{\text{\uncertain{plausible}}}{=} \mathrm{Hom}\bigl(\prod \mathcal{J}^{0}_{\mathcal{O}_D/k} / \mathrm{Ker}\, u\bigr)
\]
\note{le second argument des $\mathrm{Hom}$ n'est pas écrit}
\[
\boxed{(G\text{-}\mathcal{D}iv)/(G\text{-}\mathcal{D}iv)^{0} \subset \mathrm{Hom}\bigl(H_1(J_{*}), G(\ill{})\bigr)}
\]
\[
0 \to A' \to A \to A'' \to 0 \qquad
\boxed{\text{Ex. Néron-Severi} \subset \mathrm{Hom}\bigl(H_1(J_{*}), \mathbb{G}_m\bigr)}
\]
\note{sous $A'$, $A$, $A''$, trois petits traits obliques ; sous le second
cadre, un troisième :}
\[
\boxed{\text{A-t-on égalité ??}}
\]
\[
0 \to \prod G(\mathcal{O}_x)/G(C) \longrightarrow J^{0}_{G} \longrightarrow H^1(X, G)^{0} \to 0
\]
\[
0 \to ? \to \Bigl(\sum_D J^{*}_{k(D)/k}\Bigr)^{0} \to \mathcal{J}^{0}_{X/k} \to 0
\qquad G(L)
\]
\note{entre ces deux lignes, deux doubles traits obliques marqués « ? »,
l'un biffé ; $k(D)$ est écrit sur une lettre antérieure. À droite, un grand
arc les enferme, avec en marge : « Aucune \ill{} \ill{} \ill{} \ill{}
n'existe\ldots »}

\uncertain{Quel} était l'accouplement \uncertain{idélique} de Tate ???

\page{22}
$X \times Y$ \qquad $A$, \quad $\mathrm{Ext}^1(A, G)$

$D$ \struck{\ill{}} $G$-diviseur sur $X \times Y$,
\uncertain{horizontal} \ill{} \uncertain{fibres} de $X \times Y \to X$
\note{au-dessus de « $G$-diviseur », un mot entouré, \uncertain{ordinaire}}
transforme $0$-cycles de $X$ en $G$-diviseurs sur $Y$, et $0$-cycles du
noyau d'Albanese en $G$-diviseurs principaux (?)

\struck{\ill{}}

$X \times Y$.
\qquad \struck{\ill{}} $\to \mathrm{Pic}^{*}(Y)$
\[
J^{*}_{X/k} \longrightarrow \mathrm{Pic}^{*}_{G}(Y)
\]

\begin{tikzcd}[column sep=small, row sep=normal, nodes={font=\scriptsize}]
0 \arrow[r] & \text{Noyau d'Albanese} \arrow[r] \arrow[d] & Z^{0}_{X} \arrow[r] \arrow[d] & J^{0}_{X/k} \arrow[r] \arrow[d] & 0 \\
0 \arrow[r] & G(K)/G(Y) \arrow[r] & \mathcal{D}iv^{0}(Y, G) \arrow[r] & \mathrm{Pic}^{0}(Y, G) \arrow[r] & 0
\end{tikzcd}

\note{sous $Z^{0}_{X}$ : « $0$-cycles de degré $0$ » ; sous $J^{0}_{X/k}$ :
« Albanese de $X$ » ; sous $\mathcal{D}iv^{0}(Y,G)$ : « alg.\ équiv.\ à
$0$ », où le $Y$ est écrit sur un $X$ biffé ; sous $\mathrm{Pic}^{0}(Y,G)$ :
« classes de fibrés alg.\ équiv.\ à $0$ » ; le $G(Y)$ de la ligne du bas est
écrit sur une lettre antérieure}
\[
0 \to \text{Noyau} \longrightarrow J^{0}_{K_Y/k} \longrightarrow \mathcal{J}^{0}_{Y/k} \to 0
\]

\section{Théorème de Rosenlicht}
\note{titre inscrit de sa main en haut à droite du feuillet chamois de la
page 23, sans autre contenu}

\page{24}
\note{en haut à gauche, de sa main, « 1 » (son numéro de page ; les pages
25 à 34 ne portent pas de numéro visible)}

\emph{Théorème de Rosenlicht}

\marginal{(I), entouré}
\emph{Existence des Picard pour les courbes}

Soit $X$ un $Y$-schéma, avec
\begin{itemize}
  \item[a)] $Y$ loc.\ noethérien
  \item[b)] $X$ propre sur $Y$ (faut-il même : loc.\ projectif ?)
  \item[c)] $X$ plat sur $Y$
  \item[d)] Les fibres de $X$ sur $S$ sont des courbes algébriques
  \emph{absolument} \struck{\ill{}} et absolument \add{connexes}
  \struck{\uncertain{réduites}} [ce dernier pt n'est pas
  \uncertain{essentiel}\ldots]
  \item[e)] Le nb de composantes \struck{\ill{}} irréd.\ absolues des
  fibres de $X$ sur $Y$ est loc.\ constant sur $Y$
\end{itemize}
\note{la fin de d), à partir de « et absolument », est ajoutée en marge
droite ; « connexes » est écrit au-dessus d'un mot biffé}
\marginal{en face de e), entre crochets : si on abandonne cette hypothèse,
on \uncertain{tombe} \uncertain{sur} \uncertain{des} phénomènes à la
Néron ?!}

Alors le $S$-schéma $\mathrm{Pic}_{X/Y}$ relatif existe, et est
\emph{simple} sur $Y$.

De plus, si le nb des composantes irréductibles absolues est $n$, il y a un
revêtement étale $Y'$ de degré $n$ de $Y$ et si $\mathbb{Z}_{Y'/Y}$ est le
schéma en groupes sur $Y$ défini par $Y'$ (qui sur $Y'$ devient isomorphe à
$\mathbb{Z}^n$) on a une augmentation \struck{\uncertain{morphisme} et
\ill{}} fidèlement plate
\[
\mathrm{Pic}_{X/Y} \longrightarrow \mathbb{Z}_{Y'/Y}
\]
\marginal{\ill{}, souligné, en face de ces lignes}
\note{au-dessus de la ligne, relié à $\mathbb{Z}_{Y'/Y}$ par une accolade :
$\mathbf{N}_{Y'/Y}(\mathbb{Z}_{Y'}) = \underline{\mathrm{Hom}}_Y(Y', \mathbb{Z}_Y)$}
dont le noyau sera noté $\mathrm{Pic}^{0}_{X/Y}$, d'où la suite exacte
\[
0 \to \mathrm{Pic}^{0}_{X/Y} \to \mathrm{Pic}_{X/Y} \to \mathbb{Z}_{Y'/Y} \to 0.
\]
Notons que l'on a un \emph{homomorphisme trace}
\[
\mathbb{Z}_{Y'/Y} \longrightarrow \mathbb{Z}_Y
\]
donc une augmentation
\[
\mathrm{Pic}_{X/Y} \longrightarrow \mathbb{Z}_Y
\]
\note{cette dernière ligne et la précédente sont d'une autre encre, brune}

\page{25}
Les fibres de $\mathrm{Pic}^{0}_{X/S}$ sont connexes, et \uncertain{puisque}
\uncertain{les} fibres de $X$ sont \emph{simples}
\note{ligne ajoutée en tête de page, d'une écriture plus serrée}

Soit \struck{\ill{}} $X - S$ l'ens.\ des pts de $X$ qui \uncertain{sont}
\uncertain{non} simples sur $Y$. On va définir un morphisme
\[
\boxed{X - S \longrightarrow \mathrm{Pic}^{1}_{X/S}}
\]
où l'exposant $1$ indique qu'il est dans la composante d'augmentation $1$.
Il suffira de le \uncertain{définir} \emph{fonctoriellement} [pour le
changement de base, \uncertain{comme} \ill{}] \struck{morphisme}
\[
\Gamma(X - S/Y) \longrightarrow \mathrm{Pic}^{1}(X/S)
\]
Or on \uncertain{note} que une section de $X - S$ sur $Y$ \emph{définit un
diviseur} sur $X - S$ (transversal : $X/Y$)
\marginal{vérifier ??}
d'où un faisceau inversible, dont on constate que le \add{degré}
\struck{\ill{}} est $+1$ [car on peut supposer $Y = \mathrm{Spec}(k)$, $k$
alg.\ clos].

\marginal{(II)}
Partons d'un $X$ \emph{propre} et \emph{simple} sur $Y$, dont les fibres
sont de dim $1$, \struck{Soient $X_i$} (i.e.\ et abst.\ connexes (donc
abst.\ irréductibles)). Soit $S$ un sous-préschéma \add{fermé}
\struck{$X$ plat sur $Y$ et} ne contenant aucune fibre de $X/Y$, donc fini
sur $Y$, et plat sur $Y$. Soit $\mathfrak{J}$ le faisceau d'idéaux
définissant \uncertain{$Y$}. Alors on prouve facilement que $\mathfrak{J}$
est localement \emph{libre} [c'est purement local : un faisceau cohérent sur
$X$ plat relat.\ à $Y$ et \ill{} \ill{} une restriction \uncertain{à}
\uncertain{chaque} fibre plate sur $X$\ldots]
\note{on attendrait « définissant $S$ » ; la lettre se lit $Y$}
\marginal{en marge gauche, en face de II : fibres abst.\ connexes \ill{}
\ill{} \ill{} \ill{} \emph{\ill{}}}

\page{26}
donc \struck{il \uncertain{s'ensuit}} \add{il s'ensuit} \uncertain{facile}
que les $\mathfrak{J}^{n}$ \ill{} \ill{} $\mathcal{O}_X/\mathfrak{J}^{n}$
sont \ill{} plats sur $Y$. Soit \struck{\ill{}} $S_n$ le
sous-préschéma de $X$ défini par $\mathfrak{J}^{n}$, et soit
$X_n = X \amalg_{S_n} Y$ le \struck{\ill{}} \supplied{$Y$-}préschéma obtenu en
« \uncertain{contractant} $S_n$ en $S$ ». \struck{Soit
\uncertain{dorénavant} $X_S = X - S$.} Donc $X_0 = X$, et les $X_i$ forment
un \emph{système} \add{inductif} \struck{\uncertain{projectif}}. Donc les
$\mathrm{Pic}_{X_n/S}$ forment un système \struck{\uncertain{inductif}}
projectif, d'où un pro-($S$-préschéma en groupes) noté
$\mathrm{Pic}_{S,X/Y}$, et qui ne dépend que de la \uncertain{bête}
\emph{formée} $S$ de $X/Y$, \ill{} \ill{} \ill{} (du moins, si $Y$ est
\ill{} \uncertain{noethérien}) \ill{} résultera d'ailleurs de la
construction qu'on donnera plus bas.

\begin{tikzcd}[column sep=small, row sep=small, nodes={font=\scriptsize}]
S_n \arrow[r, "i_n"] \arrow[d, "f_n"'] & X \arrow[d, "p_n"] \\
Y \arrow[r, "\varepsilon_n"] & X_n
\end{tikzcd}

\note{à gauche de ce carré, un premier carré de même forme, hachuré et
biffé ; au-dessous : $\cdots\ X_2 \leftarrow X_1 \leftarrow X_0 = X$}

Les morphismes canoniques $X - S \to \mathrm{Pic}^{1}_{X_n/Y}$ définissent à
la limite un morphisme
\[
\boxed{\varphi_S \colon X - S \longrightarrow \mathrm{Pic}^{1}_{S, X/Y} \subset \mathrm{Pic}_{S, X/Y}}
\]

\page{27}
\emph{Théorème} Soit $G$ un $Y$-schéma en groupes \add{commutatifs} de
type fini \struck{\ill{}} sur $Y$, et $P^{1}$ un fibré principal homogène
sous $G$ [devenant trivial par extension fid.\ plate quasi-compacte\ldots],
d'où un schéma en groupes $\mathbb{Z}$-augmenté $P$, avec $P^{(1)} = P^{1}$,
$P^{(0)} = G$. Considérons l'application canonique $u \mapsto u \circ
\varphi_S$
\note{l'énoncé jusqu'ici est marqué d'un trait vertical dans la marge
gauche}
\[
\mathrm{Hom}_{\substack{Y\text{-schémas en groupes}\\ \mathbb{Z}\text{-augmentés}}}\bigl(\mathrm{Pic}_{S, X/Y}, P\bigr) \longrightarrow \mathrm{Hom}_Y(X - S, P^{1})
\]
\emph{Cette application est bijective}

(Th.\ de Rosenlicht)

\emph{Corollaire 1}
\[
\mathrm{Hom}_{Y\text{-schémas en groupes}}\bigl(\mathrm{Pic}_{S, X/Y}, G\bigr) \xrightarrow{\ \sim\ } \mathrm{Hom}_Y(X - S, G)
\]

\emph{Corollaire 2} \struck{L'\uncertain{homomorphisme} de}
\add{Considérons} \add{le} \add{foncteur} $\varphi^{*}_{S}$, de
\add{(en groupes commutatifs)} la catégorie des schémas
\emph{extensions} des $\mathrm{Pic}_{S, X/Y}$ par des $Y$-schémas en groupes
commutatifs de type fini $G$ sur $Y$, dans la catégorie
\note{« Considérons », « foncteur » et « (en groupes commutatifs) » sont
écrits dans l'interligne, au-dessus d'un passage biffé ; l'énoncé est marqué
d'un trait vertical dans la marge}
\marginal{il faut \ill{} \uncertain{définition} \ill{} « \ill{} » \ill{}
« extension » ? $G$ plat \ill{} ???}

\page{28}
des fibrés principaux homogènes sur $X - S$ de groupes de \uncertain{type}
$G$. Ce foncteur est pleinement fidèle.
\note{fin de l'énoncé du Corollaire 2, marquée d'un trait vertical en marge ;
le $X - S$ est écrit sur une lettre antérieure}

Soient $Q$, $Q'$ deux telles extensions, de groupes $G$, $G'$. Un homom.\ de
$Q'$ dans $Q$ \uncertain{est} \uncertain{formé} d'un homom.\
$u \colon G' \to G$ et d'\add{un isom.\ de $Q'' = Q \times_G G'$} \ill{}
[\struck{\ill{} isom.\ de $Q'' = Q \times_{G} G'$ avec \ill{} \ill{}
$\mathrm{Hom}_{S,u}(Q', Q)$, \ill{} \ill{} \ill{} fibré}] $Q'$
\struck{\ill{}} (\ill{} \ill{} que extensions de groupe
\emph{\uncertain{\ill{}}} $G'$). Il y a \ill{} ssi $Q'$ et $Q''$ sont
isomorphes.
\note{passage très raturé : un premier complément, biffé de plusieurs
traits, est remplacé par « un isom.\ de $Q'' = Q \times_G G'$ », écrit
au-dessus et relié par un trait}
\marginal{il faut préciser \uncertain{définition} \ill{} fibrés
\uncertain{modifiés}, attention !}

\begin{tikzcd}[column sep=small, row sep=small, nodes={font=\scriptsize}]
Q & Q'' \\
X - S \arrow[r] & \mathrm{Pic}_{S, X/S}
\end{tikzcd}

\note{croquis en marge gauche ; une flèche en pointillé monte vers $Q$, un
trait relie $Q$ à la ligne du bas}

Or
\[
\mathrm{Ext}^1\bigl(\mathrm{Pic}_{S}, G\bigr) \longrightarrow H^1(X - S, G)
\]
\note{sous le $\mathrm{Pic}$ du premier terme, un $Q$ biffé ; le $H^1$ est écrit sur un $\mathrm{Ext}$ biffé}
étant un homom.\ de groupes, pour prouver qu'il est injectif (i.e.\ $Q$ et
$Q'$ isom.\ ssi $\varphi^{*}(Q)$ et $\varphi^{*}(Q')$ le sont) on est ramené
à prouver que le noyau est nul, i.e.\ que si $\varphi^{*}(Q)$ est trivial,
alors $Q$ l'est. Cela résulte \struck{facilement} \add{aussitôt} du th.\ de
Rosenlicht.
\marginal{attention, il faut \uncertain{préciser} \ill{} \ill{}
« \ill{} » \ill{} fibrés}

Comment (\ill{} $\varphi^{*}$ est \ill{} \uncertain{pleinement} \ill{}
\ill{}) \ill{} \ill{} \ill{}, \struck{\ill{}} \ill{}

\page{29}
est ainsi ramené à prouver que
\[
\mathrm{Hom}(Q, Q) \longrightarrow \mathrm{Hom}\bigl(\varphi^{*}(Q), \varphi^{*}(Q)\bigr)
\]
est \uncertain{bijectif}. Or les deux s'identifient respect.\ à
\[
\mathrm{Hom}_{\text{gr.\ sur}\ Y}(\mathrm{Pic}_{S}, G) \quad\text{et}\quad \mathrm{Hom}_Y(X - S, G)
\]
\note{devant le second $\mathrm{Hom}$, un signe biffé}
donc est isomorphe.

En particulier

\emph{Corollaire 3} Si $G$ est un schéma en groupes commut.\ de type fini
sur $Y$, \ill{} \ill{} \ill{} homom.\ canonique \emph{injectif}
\note{l'énoncé est marqué d'un trait vertical en marge}
\[
\mathrm{Ext}^1_{\substack{\text{schémas en groupes}\\ \text{sur}\ Y}}\bigl(\mathrm{Pic}_{S, X/S}, G\bigr) \longrightarrow H^1(X - S, G)
\]
\note{l'$\mathrm{Ext}$ est écrit sur un $H$ biffé}
\marginal{préciser la signification\ldots}

Il \emph{semble} que \uncertain{c'est} \emph{bijectif} \uncertain{chaque}
fois que $S \to Y$ est \emph{surjectif}, i.e.\ que les fibres de $X/Y$ sont
non propres, \ill{} \emph{\ill{}}. Dans le cas où $S = \emptyset$,
$G = \mathbb{G}_m$, \ill{} \ill{} \ill{} pas surjectif.
\note{« bijectif » est écrit sur un mot antérieur}

\page{30}
\note{page de prose rapide, sans formule centrée ; seuls les termes
mathématiques et quelques mots se lisent}
Je \uncertain{remarque} néanmoins que dans le cas général, pour prouver que
un fibré principal $E$ sur $X - S$ provient d'une \struck{fibré principal}
extension de $\mathrm{Pic}_S$ par $G$, il suffit (utilisant le théorème de
la descente) de le vérifier en \uncertain{remplaçant} $S$ par les spectres
des \emph{complétés} des \ill{} anneaux locaux de \ill{} [\ill{} plus,
utilisant la géométrie formelle \struck{\ill{}} \ill{}]. $G$ \ill{} bien
\ill{} \ill{} \ill{} \ill{} \ill{} \ill{} \emph{formellement}. \ill{} \ill{}
\ill{} \ill{} \ill{} d'un \ill{} \ill{} \ill{}. Travaillant de proche en
proche, on est \ill{} à \uncertain{vérifier} que \ill{} \ill{} \ill{}
\add{\ill{}} les corps résiduels de $S$. Utilisant \ill{} \ill{} \ill{}
\uncertain{suivant} on est ramené à le vérifier \ill{} \ill{} les
clôtures algébriques des corps résiduels\ldots\ Ce point \ill{} \ill{}
\ill{} \ill{} d'un tiers plus \ill{}ment.

\page{31}
Utilisant une technique \uncertain{de} \ill{} \ill{}, \uncertain{il}
\uncertain{semble} que, on \struck{et} définit \ill{} \ill{} un
homomorphisme
\[
\longleftrightarrow\ H^1(X, G) \xrightarrow{\ d\ } \mathrm{Hom}_{Y\text{-gr}}(\mathbb{G}_m, G)
\qquad \text{[\ill{} \emph{degré}]}
\]
\note{devant la double flèche, une lettre biffée, \uncertain{$Y$}}
\marginal{en oblique dans la marge gauche, reliée par la double flèche :
\uncertain{la} définition \ill{} \ill{} \ill{} \ill{} \uncertain{hyp.},
\ill{} $= \mathrm{Spec}(\ill{})$, $G = \mathbb{G}_m$ \ill{}}
et on trouve que la suite
\[
\boxed{0 \to \mathrm{Ext}^1(\mathrm{Pic}_{X/Y}, G) \to H^1(X, G) \to \mathrm{Hom}_{Y\text{-gr}}(\mathbb{G}_m, G)}
\]
\note{la barre du premier $\mathrm{Ext}$ est un $\mathrm{Pic}$ surchargé ;
sous le second $H^1$, l'argument se lit $X$, suivi d'un indice
\ill{}}
est exacte. Je \uncertain{ne} \uncertain{sais} \uncertain{pas} \ill{}
\ill{} \ill{} la suite exacte \uncertain{plus} \uncertain{complète} :
\[
\boxed{\to \mathrm{Ext}^2(\mathrm{Pic}_{X/Y}, G) \to H^2(X, G) \to \mathrm{Ext}^1_{Y\text{-gr}}(\mathbb{G}_m, G) \to \cdots}
\]
\ill{} a \ill{} \ill{} \ill{}
\[
\mathrm{Ext}^1(\mathbb{G}_m, G) \longrightarrow \mathrm{Ext}^{\ill{}}(\mathrm{Pic}_{X/Y}, G)
\]
\ill{} la multiplication par un élément \emph{fondamental}
\[
\xi \in \mathrm{Ext}^2_{Y\text{-groupes}}(\mathrm{Pic}_{X/Y}, \mathbb{G}_m)
\]
dont la définition \ill{} \ill{} \ill{} \ill{} ; il \ill{} \ill{}
\ill{} [\uncertain{cela} s'il existe \ill{} $S \to Y$ non
\uncertain{surjectif}],

\page{32}
\uncertain{condition} toujours \uncertain{remplie} lorsque $S$ est
\uncertain{local} : l'on considère l'extension
\[
0 \to N_S \longrightarrow \mathrm{Pic}_{S, X/Y} \longrightarrow \mathrm{Pic}_{X/Y} \to 0
\]
d'où un
\[
\xi \in \mathrm{Ext}^1(\mathrm{Pic}_{X/Y}, N_S)
\]
puis l'extension
\[
0 \to \mathbb{G}_{m,Y} \longrightarrow \mathcal{E}_{X/Y/S} \longrightarrow N_S \to 0
\]
d'où un
\[
\eta \in \mathrm{Ext}^1(N_S, \mathbb{G}_{m,Y})
\]
et on pose
\[
\zeta = \eta \circ \xi
\]
\emph{Dans le cas où $\zeta = 0$}, \emph{i.e.\ par exemple s'il existe une
section de $X$ sur $Y$} (qui \ill{} \ill{} \ill{} \ill{} $\zeta$, \ill{}
\ill{} \ill{} \ill{} \ill{} $\eta = 0$] \ill{} \ill{} \ill{} \ill{} les
\uncertain{suites} exactes \uncertain{plausibles} \uncertain{suivantes}
\[
\boxed{0 \to \mathrm{Ext}^{i}(\mathrm{Pic}_{X/Y}, G) \to H^{i}(X, G) \to \mathrm{Ext}^{i-1}(\mathbb{G}_m, G) \to 0}
\]
\note{la lettre notée $\mathcal{E}$ est une majuscule ornée, écrite sur un signe biffé}

\page{33}
\note{feuillet plus large, en deux colonnes ; la colonne de gauche est
donnée d'abord}
Passant à la limite sur des $S$ de plus en plus \uncertain{grands} (cet
\ill{} \uncertain{filtrant} sur $Y$, \ill{} \ill{} la \ill{}) on trouve un
\ill{} \add{pro-$Y$-}\ill{}
\[
\mathrm{Pic}_{\infty\, X/Y},
\]
\struck{\ill{}} et un \struck{\ill{}} \add{pro-morphisme} \ill{} \ill{}
fibre de $X/Y$
\[
\varphi_\infty \colon X \longrightarrow \mathrm{Pic}^{1}_{\infty}
\]
qui est solution d'un problème universel relatif aux pro-morphismes
\ill{} \ill{} fibres de $X/Y$ dans les pro-\ill{} \add{commutatifs} des
$Y$-schémas en groupes (de type fini. \ill{} \ill{}, si $R_{X/Y}$
\struck{\ill{}} désigne \ill{} \emph{\ill{} \ill{}} sur $Y$ défini par $X$
\[
\left\lbrace
\begin{array}{l}
\mathrm{Hom}(R_{X/Y}, P^{1}) \simeq \mathrm{Hom}_{Y\text{-gr}}(\mathrm{Pic}^{1}_{\infty}, P) \\
\underline{H}^{0}(R_{X/Y}, G) \xleftarrow{\ \sim\ } \mathrm{Hom}_{Y\text{-gr}}(\mathrm{Pic}_{\infty}, G) \\
\underline{H}^{1}(R_{X/Y}, G) \xleftarrow{\ \sim\ } \mathrm{Ext}^{1}_{Y\text{-gr}}(\mathrm{Pic}_{\infty}, G) \quad (\text{et } \ill{}\ \ill{}) \\
\underline{H}^{i}(R_{X/Y}, G) \xleftarrow{\ \sim\ } \mathrm{Ext}^{i}_{Y\text{-gr}}(\mathrm{Pic}_{\infty}, G)
\end{array}\right.
\]
\note{devant l'accolade, un « $H^1(Z,$ » biffé ; en marge, reliées à
l'accolade : « \ill{} des \uncertain{faisceaux} sur $S$ » ; les
$\underline{H}$ sont soulignés d'un trait ondulé ; dans le premier
$\mathrm{Hom}$, un indice \ill{} sous $R$}

Je \uncertain{note} : \uncertain{liés} les
$\mathrm{Ext}^{i}(\mathrm{Pic}_{S, X/S}, G)$ avec les
$\mathrm{Ext}^{i}(\mathrm{Pic}^{0}_{S, X/S}, G)$ \uncertain{par}
\uncertain{l'intermédiaire} d'habitude. Le lien se fait par la suite
exacte des $\mathrm{Ext}$ \ill{} de
\[
(*) \qquad 0 \to \mathrm{Pic}^{0}_{S} \to \mathrm{Pic}_{S} \to \mathbb{Z} \to 0
\]
donnant une
\[
\begin{aligned}
&\to H^0(S, G) \to \mathrm{Hom}(\mathrm{Pic}_S, G) \to \mathrm{Hom}(\mathrm{Pic}^{0}_{S}, G) \to H^1(S, G) \to \mathrm{Ext}^1(\mathrm{Pic}_S, G) \\
&\to \mathrm{Ext}^1(\mathrm{Pic}^{0}_{S}, G) \to H^2(S, G) \to \mathrm{Ext}^2(\mathrm{Pic}_S, G) \to \cdots
\end{aligned}
\]
\note{au-dessus de $\mathrm{Hom}(\mathrm{Pic}_S, G)$, relié par $\simeq$ :
$\mathrm{Hom}(X - S, G)$ ; au-dessus de $\mathrm{Ext}^1(\mathrm{Pic}_S, G)$ :
$H^1(X - S, G)$, relié par « $\simeq$ ? » ; au-dessus de
$\mathrm{Ext}^2(\mathrm{Pic}_S, G)$ : $H^2(X - S, G)$, relié par
« $\simeq$ ? » ; au-dessous de la suite, un grand signe biffé}

On \uncertain{notera} que dans cette théorie, [du moins pour les
\ill{}] on n'a aucun besoin de notion de \ill{} \uncertain{local}
\ill{} aucune \uncertain{notion} de la théorie des Jacobiennes locales, qui
\ill{} \ill{} \ill{} \ill{} \ill{}\ldots

\page{34}
\uncertain{Détailler}
\[
0 \to \underset{\substack{\text{faisceau de}\\ \text{Picard relatif}\\ \text{\uncertain{ordinaire}\ (?)}}}{\mathrm{Pic}(X/Y)_0} \longrightarrow \underline{\mathrm{Ext}}^1\bigl(\mathrm{Pic}^{0}_{X/Y}, \mathbb{G}_m\bigr)
\]
\note{« ordinaire » est souligné ; sous $\underline{\mathrm{Ext}}^1$, le
soulignement est ondulé}

\begin{tikzcd}[column sep=small, row sep=normal, nodes={font=\scriptsize}]
\underline{\mathrm{Ext}}^1(\mathrm{Pic}^{0}_{X/Y}, \mathbb{G}_m) \arrow[d, "2"] & \text{\emph{autodualité des jacobiennes}} \\
\underline{H}^{0}(Y, \mathrm{Pic}^{0}_{X/Y}) &
\end{tikzcd}

\note{la flèche verticale, marquée « 2 », descend de l'$\underline{\mathrm{Ext}}^1$
de la ligne précédente ; devant $\underline{H}^{0}$, un signe biffé}

La multiplicité \ill{} \ill{} \ill{} par les faisceaux de
\uncertain{Weil} [\uncertain{sur} $Y$]. Ceci dit, on \ill{} avoir \ill{} des
\ill{} \uncertain{convenables}
\[
\underline{\underline{H}}^{*}(X_S/Y, G) \xleftarrow{\ \sim\ } \underline{\underline{\mathrm{Ext}}}^{*}_{Y}(\mathrm{Pic}_S, G)
\]
\note{sous le second terme, une flèche marquée « 2 » descend vers
$\underline{\underline{\mathrm{Ext}}}^{*}(\mathrm{Pic}^{0}_{S}, G)$, suivi de
\ill{} \ill{} doublement soulignés ; les soulignements doubles sont les
siens}
D'autre part on a \add{des} suites spectrales,
\[
\begin{aligned}
\mathrm{Ext}^{*}(\mathrm{Pic}_S, G) &\Longleftarrow H^{p}\bigl(Y, \underline{\underline{\mathrm{Ext}}}^{q}(\mathrm{Pic}_S, G)\bigr) \\
\underline{H}^{*}(X_S/Y, G) &\Longleftarrow H^{p}\bigl(Y, \underline{\underline{H}}^{q}(X_S/Y, G)\bigr)
\end{aligned}
\]
d'où \ill{} \struck{substituant une suite spectrale}
\struck{$H^{*}(X_S/Y, G) \Longleftarrow H^{p}\bigl(Y, \underline{\underline{\mathrm{Ext}}}^{q}(\mathrm{Pic}_S, G)\bigr)$}
\struck{d'où une \ill{} \ill{} \ill{} \ill{}}
\note{la formule encadrée et la ligne qui la suit sont barrées d'un trait
en zigzag et d'un trait oblique}
\uncertain{compatible} \ill{} résultat plausible
\[
H^{*}(X_S/Y, G) \simeq \mathrm{Ext}^{*}_{S}(\mathrm{Pic}_S, G)
\]
\marginal{\struck{$0 \to H^1(Y, \underline{\mathrm{Hom}}(\mathrm{Pic}_S, G))$}}

\section{Adèles des préschémas}
\note{titre inscrit de sa main en haut à droite du feuillet chamois de la
page 35, sans autre contenu}

\page{36}
$\mathbb{S}_\infty$ schéma d'anneaux \add{affine} sur $\mathbb{Z}$, si $A$
est un anneau, on a
\[
\mathbb{S}_\infty(A) = A[[t]]^{**}
\]
\note{l'argument de $\mathbb{S}_\infty$ est écrit sur une lettre biffée ;
$\mathbb{S}$ rend son $S$ barré, à la manière d'une lettre double}
(séries formelles d'augmentation $1$) où l'addition est la multiplication
des séries formelles, et le produit est la convolution $*$. Il y a aussi
des opérations $\lambda^{i}$. \uncertain{Rappelons},
\[
(1 + at) * f(t) = f(at) \qquad \text{i.e.} \qquad \{1 + at\}\{f(t)\} = \{f(at)\}
\]
\marginal{à l'encre rouge, en oblique dans la marge gauche : Relations entre
séries formelles et vecteurs de Witt}
Si $A$ est de car $p > 0$, on aura un \emph{endomorphisme} $F$ de
$\mathbb{S}_\infty(A)$ [provenant de l'endomorphisme de Frobenius de
$\mathbb{S}_\infty\, \mathbb{Z}/p\mathbb{Z}$]
\[
F\Bigl(\sum a_i t^{i}\Bigr) = \sum a_i^{p} t^{i}
\]
et on aura
\[
f(t)^{p} = F(f)(t^{p})
\]
\note{devant cette formule, un début biffé, « $F(f(t))$ »}
Cela conduit alors à poser
\[
V(f)(t) = f(t^{p})
\]
[$V$ est un endomorphisme additif] et on aura
\[
FV(\{f\}) = VF(\{f\}) = p\{f\}
\]
N.B.\ $F$ est aussi un endomorphisme pour la structure d'anneau.

\struck{\ill{} \ill{}} par exemple
\[
(1 + x t^{p^{i}}) * f(t) = F^{i}(f)(x t^{p^{i}}) \qquad\quad (1 + at)^{p^{i}} * f = f(at)^{p^{i}}
\]
car \struck{si $x = a^{p^{i}}$, le premier membre est $1 + x$}
\[
= F^{i}(f)(a^{p^{i}} t^{p^{i}}) = F^{i}(f)(x t^{p^{i}}).
\]
\note{dans la première formule, le $x$ est écrit sur un $a$}

Sans restriction \struck{\ill{}} de caractéristique, on a un morphisme
d'anneaux
\[
\psi \colon \mathbb{S}_\infty \longrightarrow \mathbb{W}_\infty
\qquad \text{[anneau de Witt pour le nb premier } p\text{]}
\]
\marginal{« ? » en marge gauche}
qui \ill{} \ill{} \ill{} morphisme additif \uncertain{déterminé} par
\[
\psi(\{1 + at\}) = \varepsilon(a) = (a, 0, 0, \ldots)
\]

\page{37}
(cela montre \ill{} \ill{} \ill{} l'application des « représentants
multiplicatifs » est fonctorielle, i.e.\ provient d'un morphisme de
$\mathbb{Z}$-schémas $\mathbb{G}_m \to \mathbb{W}_\infty$).

\page{39}
\[
0 \to K^{*} \longrightarrow \prod_{\text{local}} K_x^{*} \longrightarrow J_C^{(\bar k)} \to 0
\]
\marginal{un grand « ? » entouré, en haut à droite}
\marginal{dans un cadre arrondi, à gauche : \uncertain{Énoncé} de
\uncertain{passage} à la limite}
\[
\begin{array}{ll}
\hat H^{*}(\Gamma, K^{*}) = 0 & \text{par le théorème de \ill{} global} \\
\hat H^{*}(\Gamma, K_x^{*}) = 0 & \text{\quad " \qquad local}
\end{array}
\]
d'où
\[
\hat H^{*}(\Gamma, J_C^{(\bar k)}) = 0 \qquad \text{\uncertain{géométriques}}
\]
En particulier
\struck{$C^{*}(\Gamma, J_C(\ill{}))$ est une résolution de $J$.
$H^0(C^{*}(\Gamma, J_C(\bar k)))$}
\[
C^{*}(\Gamma, J_C^{(\bar k)}) \quad \text{est une \emph{résolution} de} \quad H^0(\Gamma, J_C^{(\bar k)})
\]
Or \ill{} \ill{}, on sait que $H^1(\Gamma, K^{*}) = 0$
\[
\begin{aligned}
H^0(\Gamma, J_C(\bar k)) &= H^0\Bigl(\prod_{\substack{\text{local}\\ x \in X}} K_x^{*}\Bigr) \Big/ \mathrm{Im}\, H^0(K^{*}) \\
&= \Bigl(\prod_{\substack{\text{local}\\ x' \in X'}} K'^{*}_{x'}\Bigr) \Big/ K'^{*} = J_{C'}(\bar k)
\end{aligned}
\]
C'est bien une résolution (Mais qui ne \uncertain{sont} pas
« admissibles ») de $\pi$-\uncertain{groupes}

D'où
\[
H^{*}(\pi, J_{C'}(\bar k)) \Longleftarrow H^{p}\bigl(H^{q}(\pi, \underbrace{C^{*}(\Gamma, J_C(\bar k))}_{C^{*}(\Gamma, J_C)(\bar k)})\bigr)
\]
\[
q > 0 \qquad H^{q}(\pi, C^{*}(\Gamma, J_C)(\bar k)) = H^{q}(\pi, C^{*}(\Gamma, \mathbb{Z})) = 0 \quad (\text{Lang})
\]
\note{au-dessus du second membre, « $C^{*}C$ » biffé}
\[
q = 0 \qquad C^{*}(\Gamma, J_C(\bar k))^{\pi} = C^{*}(\Gamma, J_C(k))
\]
\note{en marge droite : « trivial car $C^{*}(\Gamma, \mathbb{Z}) \simeq
\mathbb{Z}^{(\Gamma)}$ », suivi d'un crochet et d'un trait ondulé}

\page{40}
\note{page écrite tête-bêche par rapport à la numérotation des archivistes ;
formules éparses, données de haut en bas dans le sens de l'écriture}
\[
H^{p}\bigl(k, \underline{\underline{H}}^{q}(\Gamma, J_{C'})\bigr) \Longleftarrow H^{p}(k, J_C) \qquad J_C(k)
\]
\note{entre les deux membres, un signe biffé ; devant $J_C(k)$, un signe
biffé d'un trait oblique}
\[
H^{p}\bigl(H^{q}(k, \underline{\underline{C}}^{*}(\Gamma, J_{C'}))\bigr)
\]
\[
\underset{q = 0}{H^{p}}\bigl(C^{*}(\Gamma, J_{C'})\bigr) = H^{p}(\Gamma, J_{C'})
\]
\[
J_{C'}(\bar k) \qquad \underline{C}(\Gamma^{n}, J_{C'})(\bar k) \qquad \text{\emph{résolution} de } J_C(\bar k) \text{ \uncertain{par} des } \pi\text{-\uncertain{modules}}
\]
$\Gamma$ \quad $\pi$
\[
H^{*}(\pi, J_C(\bar k)) \Longleftarrow H^{p}\bigl(H^{q}(\pi, \underline{C}(\Gamma^{\cdot}, J_{C'})(\bar k))\bigr)
\]
\[
\begin{array}{ll}
\| & \\
J_C(k) & H^{q}(\pi, \underline{C}(\Gamma^{\cdot}, \mathbb{Z})) \underset{q > 0}{=} 0
\end{array}
\]
\[
H^0(\Gamma, J_{C'}(\bar k)) \qquad C(\Gamma^{\cdot}, J_{C'}) \qquad H^{p}(\Gamma, J_{C'}(k))
\]
\note{devant $C(\Gamma^{\cdot}, J_{C'})$, un « $H^0($ » biffé}
\[
\underline{C} \qquad \Gamma \qquad C(\Gamma^{n}, J^{\cdot}_{C'})
\]
\[
J_C(\bar k) \to \underline{J_{C'}(\bar k)} \longrightarrow J_{C'}(\bar k)^{\Gamma} \longrightarrow J_{C'}(\bar k)^{\Gamma^{2}}
\]
\[
(\delta\varphi)(\underset{\uparrow\,\gamma}{x}) = \varphi(\gamma x) - \varphi(x) \qquad \tau_\gamma \varphi - \varphi
\]

\end{document}
